% Slide 2 – problem and target group (criterion: relevance and usefulness for the group, 25 %).
\begin{frame}{Reliable information instead of guesswork}
  \begin{columns}[T,onlytextwidth]
    \begin{column}{0.49\textwidth}
      \begin{block}{Problem}
        \begin{itemize}
          \item Accessibility data is scattered, incomplete and quickly goes out of date.
          \item Without a source and a date, users cannot tell whether to trust it.
          \item Keeping a complete database for a whole city up to date by hand does not scale.
        \end{itemize}
      \end{block}
    \end{column}
    \begin{column}{0.49\textwidth}
      \begin{block}{Who it is for}
        People with access needs and their carers, parents with prams, older people, people recovering from
        medical procedures, tourists with luggage or pets, and the medical tourism industry.
      \end{block}
    \end{column}
  \end{columns}
  \vspace{1mm}
  \begin{exampleblock}{Accessly}
    Combines public data, information from venue owners and confirmations from the community.\\
    Every piece of information has a source, a date and a verification status.
  \end{exampleblock}

  \note[item]{Open with the story: Anna, a wheelchair user, wants a coffee in the centre. Last year's map says “accessible”; today there is a step at the door and nobody knows whether the toilet is downstairs.}
  \note[item]{Primary group from the brief: wheelchairs and prams (the “Wheelchair” and “Pram” profiles). The other profiles reuse the same attributes, so they cost nothing extra.}
  \note[item]{Stress the two groups: end users and owners; without owners the data will not stay current.}
  \note[item]{A profile is a set of preferences (steps, lift, quiet), not disability information; this comes back on the privacy slide.}
\end{frame}
